\begin{textsample}{POS Dim 1 – vem\_ed\_15 – Score 22.00 – vem\_ed\_15/t065.txt}  \label{ex:f1_pos_097}
PERSPECTIVAINTERNACIONAL A \textbf{parte} da manhã (9h30 às 12h) \textbf{foi} dedicada aos relatos de parceiros internacionais das Fatecs na realização de PCIs.
Da China, as professoras Sally Xu e Hang Yin (Tianjin Normal University) \textbf{compartilharam} a experiência de um PCI desenvolvido há quatro anos com Fatecs, com uma participação de aproximadamente 1.000 alunos brasileiros e chineses por ano.
Voltado às \textbf{competências} interculturais e ao \textbf{aprendizado} do inglês por \textbf{parte} de falantes não-nativos, envolveu nove unidades em 2022: Campinas, Catanduva, Franca, Ferraz de Vasconcelos, Itaquaquecetuba, Mogi das Cruzes, Piracicaba, São José do Rio Preto e Sebrae.
Adya Sharma, diretora do Symbiosis Centre for Management Studies (Symbiosis International, Deemed University)/Índia, apresentou o país, a instituição e os PCIs / COIL \textbf{desenvolvidos} com Fatecs.
Os professores Sonica Rautela, Nehajoan Panackal e Ashutosh Mathur comentaram o \textbf{desenvolvimento} de um estudo comparativo na área de marketing.
Estudantes da IES indiana também gravaram depoimentos, celebrando a aprendizagem sobre novas culturas, a \textbf{interação} e a ampliação de \textbf{horizontes}.
Divinia Jithoo, especialista em educação internacional na Durban University of Technology (DUT / África do Sul), comentou sobre a forte parceria entre DUT e as Fatecs.
Como \textbf{são} duas instituições no eixo de cooperação Sul-Sul, os estudantes enfrentam desafios similares e, devido às limitações financeiras, a realização de \textbf{intercâmbios} virtuais \textbf{torna} as \textbf{ações} pedagógicas mais viáveis.
Lindewa Mkhize, coordenadora de educação internacional da DUT, participou ao vivo da sessão de perguntas e \textbf{destacou} a colaboração entre coordenadores, professores e alunos.
Hope Windle, diretora do SUNY COIL Center (EUA), \textbf{contou} como a iniciativa COIL se \textbf{expandiu} pelo mundo, incluindo as Fatecs do Centro Paula Souza.
Trata-se de um relacionamento de \textbf{longa} data, \textbf{que} \textbf{iniciou} com um projeto realizado em 2013 entre SUNY Ulster e Fatec Americana.
Entre as principais lições aprendidas estão: ouvir os parceiros, \textbf{ser} flexível, definir os \textbf{papéis} de \textbf{cada} um e refletir sobre a experiência.
Javier Guerrero, coordenador de \textbf{desenvolvimento} de professores da Uniminuto (Colômbia), comentou sobre os projetos realizados desde 2021 com Fatecs, envolvendo áreas de administração e licenciaturas em inglês e espanhol.
Ressaltou a transferência de \textbf{conhecimento} pela equipe dos PCIs na realização de \textbf{intercâmbios} virtuais e capacitação de professores da Uniminuto.
A professora Bianca Cely \textbf{mencionou} a \textbf{importância} da realização de eventos acadêmicos em conjunto como o 1ºColóquio de Literatura Brasileira, realizado em 2021com a participação de Regiane Camargo (equipe dos PCIs / Cesu) e Wilton Garcia (Fatec Itaquaquecetuba).

% matched lemmas: aprendizado, ação, cada, compartilhar, competência, conhecimento, contar, desenvolvimento, desenvolvir, destacar, expandir, horizonte, importância, iniciar, interação, intercâmbio, longo, mencionar, papel, parte, que, ser, tornar
\end{textsample}
